\begin{abstract}
A strongly lensed gravitational-wave (GW) source produces several images that arrive minutes to months apart, and galaxies capable of lensing it are already catalogued in their tens of thousands. We search all public LIGO--Virgo--KAGRA data for pairs of events that are two images of one source lensed by a known galaxy. The first stage ranks every pair among 842 events and candidates from the first through fourth observing runs (354,061 pairs) by the consistency of their masses, sky positions, arrival delay and apparent distances with each of 47,227 catalogued galaxy lenses; the chance level of the ranking statistic is set by rotating the lens catalogue on the sky. The second stage analyses the 56 best pairs with a joint Bayesian fit of both signals to the public strain, giving the coherence ratio between the lensed and the unrelated-event hypotheses. Its null distribution is measured on unrelated simulated events in real detector noise, analysed identically and paired only when they pass the same stage-1 selection as the real pairs; pair-matched backgrounds without this selection underestimate it by several units. No pair is significant: the largest stage-1 coincidence is reproduced by chance in 57\% of catalogue rotations, and the largest coherence ratio of a pair with clean data, $\lnC=13.4$ for GW250109\_010541 and GW250116\_015318, has a false-alarm probability of 0.92 after accounting for the 42 pairs tested. Simulated lensed pairs from the GWTC-4.0 population injected into the same noise are selected for stage~2 in 82\% of cases, and 43\% of the selected ones exceed the largest background coherence ratio. This bounds the number of lensed pairs with both images above SNR 10 in O4a and O4b to fewer than 10 at 90\% confidence. We also show that unmitigated non-Gaussian transients in the public strain produce the two largest coherence ratios of the search, 46 and 26, and give a test that identifies them from the maximum likelihood alone.
\end{abstract}
